\section{Common paired-cap sewing}\label{sec:sewing}

Let \(G=(V,E)\) be a finite simple connected graph with at least one
edge and degree \(d_v\) at \(v\). The unfilled vertex block is
\(X_v=\Sigma_{0,d_v+2}\times S^1\) with base-then-fibre orientation.
At every port \(z\) is the oriented base section and \(h\) is the
fibre, with \(z\cdot h=+1\). Fill two specified ports by the primitive
slopes \(p_vz+h\) and \(-p_vz+h\), with \(p_v\in\ZZ_{>0}\).
On the other ports, glue the two ends of edge \(e=vw\) by
\[
        z_{ve}\longmapsto \epsilon_e h_{we},\qquad
        h_v\longmapsto \epsilon_e z_{we},\qquad \epsilon_e\in\{+1,-1\}.
\]
Both maps reverse boundary orientation. Write \(M_{G,\epsilon,\mathbf p}\)
for this closed oriented manifold. Default cap orders are \(p_v=2\), and
then we suppress \(\mathbf p\). Write \(A_\epsilon\) for the signed
adjacency, \(n=|V|\), \(e=|E|\), and \(b_G=e-n+1\).

\subsection{Composition phase and the sign of a seam}

In the torus spaces use the section Lagrangian and the core basis
paired against its oriented dual. The basic sewing formula is the
Verlinde product-block contraction with the modular cap and Hopf seam
states \cite{BHMV,P50,P50B}. For a zero-weight lift preserving the
section Lagrangian, reversal of both torus coordinates is the
involution \(C=S^2\), charge conjugation. On an unpunctured torus this
can also be read by reversing the solid-torus core in the dual-paired
basis. Thus the change of sign of a seam inserts the identity in any
theory with self-dual simple labels. This statement is about the bare
contractions; it does not by itself determine their framing phase.

Here is the global phase check, including varying cap orders. Let
\(q,p\) be the section and fibre coordinates on the combined port
space, and let \(U\) take their per-vertex sums. On the block side the
homology kernel is
\(L_X=\{(U^{\mathsf T}a,p):Up=0\}\); on the paired cap and seam side
it is \(L_Y=\{(Tp,p)\}\). For the latter description, \(T\) has the cap
entries \(p_v,-p_v\) and the seam blocks
\(\left(\begin{smallmatrix}0&-\epsilon_e\\-\epsilon_e&0\end{smallmatrix}\right)\).
The two reciprocal cap entries cancel in the sum at each vertex.
Thus \(T\) is invertible, \(\sig T=0\), and \(UT^{-1}U^{\mathsf T}=-A_\epsilon\).
With the section reference \(L_0=\{(q,0)\}\), the cap/seam-to-block
Maslov form is the restriction of \(T\) to \(\ker U\)
\cite{GM,P50B}. This follows from writing the representative
\((Tp,p)\) as \((0,p)+(Tp,0)\) for \(Up=0\). Its signature is
\[
 \sig(T|_{\ker U})=\sig T-\sig(UT^{-1}U^{\mathsf T})=\sig A_\epsilon .
\]
For the degenerate case, the \(T\)-orthogonal complement of \(\ker U\)
is \(T^{-1}\operatorname{im}U^{\mathsf T}\). A joint radical of the two
restrictions and a dual isotropic space split off as a hyperbolic
summand of signature zero. This gives the same signature formula
without an invertibility assumption on \(A_\epsilon\).
On a bipartite graph \(\diag(+1,-1)\), with signs constant on the
parts, conjugates \(A_\epsilon\) to \(-A_\epsilon\), so the signature
and the entire composition anomaly vanish. In particular this
argument proves equality of \emph{oriented complex} scalars, not only
of absolute values. Interior ribbon graphs leave these boundary
homology kernels unchanged.

\subsection{Exact \texorpdfstring{\(\SU(2)\)}{SU(2)} weights for cap order two}\label{sec:weights}

Use shifted level \(r=k+2\ge3\) and
\[
 Z_r(S^2\times S^1)=1,\qquad Z_r(S^3)=S_{11},\qquad
 \tau_r(M)=Z_r(M)/S_{11}.
\]
For core labels \(\ell,j=1,\ldots,r-1\) the modular matrices and
anomaly unit are
\[
 S_{\ell j}=\sqrt{\frac2r}\sin\frac{\pi\ell j}{r},\qquad
 \Theta_{\ell\ell}=e^{\pi i(\ell^2-1)/(2r)},\qquad
 \kappa_r=e^{3\pi i/4-3\pi i/(2r)}.
\]
At vertex \(v\) the paired cap has weight
\(\big|(S\Theta^2 S)_{\ell1}\big|^2\) and the remaining local weight
is \(S_{1\ell}^{-d_v}\). A seam contributes \(S_{\ell j}\).
The two cap states may exchange roles according to the geometric
framing convention without changing their product. The closed
contraction includes a trace along a cycle, with no divide by its
state-space dimension.

For clarity, the parity in the cap weight can be checked directly.
Let \(G_r(j)=\sum_{u\bmod 2r}e^{\pi i(u^2+ju)/r}\). The cosine
difference formula for the sine product gives
\((S\Theta^2 S)_{\ell1}\), up to a unit phase, as
\([G_r(\ell-1)-G_r(\ell+1)]/(2r)\).
Translation gives \(G_r(j+2)=e^{-\pi i(j+1)/r}G_r(j)\).
In the squared modulus of the complete sum, Fourier orthogonality
retains only the two differences \(0,r\). Hence
\(\left|G_r(j)\right|^2=4r\) if \(r+j\) is even and is zero
otherwise. With the defect label \(a=r-\ell\), the cap weight
therefore equals \(4r^{-1}\cos^2(\pi a/(2r))\) on odd \(a\) and
vanishes for even \(a\).

Put
\[
 K(x)=\frac{\sin(\pi x)}{\pi x},\quad K(0)=1,\qquad
 F_d(x)=\cos^2(\pi x/2)\left(\frac{\pi x}{\sin\pi x}\right)^d .
\]
Converting to odd defects gives the exact zero-insertion formula
\begin{align}
 Z_r(M_{G,\epsilon})
 &=(-1)^{re}\kappa_r^{\sig A_\epsilon}
   2^{2n-e/2}\pi^{-e}r^{3e/2-n}I_r(G),\label{eq:finite}\\
 I_r(G)
 &= \sum_{\substack{1\le a_v<r\\ a_v\ {\rm odd}}}
       \prod_{v\in V}F_{d_v}(a_v/r)
       \prod_{vw\in E}K(a_v a_w/r).\notag
\end{align}
In fact the two odd defects at the ends of an edge give
\(\sin(\pi(r-a_v)(r-a_w)/r)=(-1)^r\sin(\pi a_va_w/r)\).
The ensuing \(a_v^{d_v}\pi^{d_v}/r^{d_v}\) precisely cancels the
corresponding factors introduced by \(F_{d_v}\).
The remaining powers follow from \(n\) paired caps \(4/r\),
\(e\) Hopf factors \(\sqrt{2/r}\), and \(2e\) occurrences of
\(S_{1a}^{-1}\) in the vertex powers. The complete oriented formula
agrees with \cite{P50,P50B}. The sign \((-1)^{re}\) appears from
the defect coordinates, not the choice of seam signing.
